\section*{35. Über Maximalbereiche aus ganzzahligen Funktionen.}

\begin{center}
Von Emmy Noether (Göttingen).\par
\emph{Rec. Soc. Math. Moscou 36 (1929), S. 65--72}
\end{center}




Die folgenden Untersuchungen ber"uhren sich nahe mit dem Hilbert'schen Problem der relativ-ganzen Funktionen\textsuperscript{1)}; sie gehen auf eine gelegentliche Frage von E.\ Hecke zur"uck, der einen Spezialfall --- den er als Hilfssatz brauchte --- direkt rechnerisch erledigen konnte\textsuperscript{2)}.

Die Fragestellung ist die folgende: Unter einem \textit{Maximalbereich} aus ganzzahligen Funktionen in $x_1,\ldots,x_n$ --- Polynomen, Potenzreihen oder Quotienten solcher --- innerhalb eines gegebenen Bereiches $\mathfrak{Z}$ verstehe ich einen solchen, der sich durch Zuf"ugung von ganzzahligen Nennern nicht mehr erweitern l"a\ss t, bei dem also aus der Zugeh"origkeit von $a \cdot g(x)$ --- wo $g(x)$ in $\mathfrak{Z}$ --- stets die Zugeh"origkeit von $g(x)$ folgt, unter $a$ eine von Null verschiedene ganze rationale Zahl verstanden. Durch jeden Integrit"atsbereich $\mathfrak{I}$ in $\mathfrak{Z}$ wird eindeutig ein maximaler $\mathfrak{M}(\mathfrak{I})$ erzeugt, der aus allen Funktionen $g(x)$ aus $\mathfrak{Z}$ besteht, f"ur die es mindestens ein $a \neq 0$ gibt, so da\ss{} $a \cdot g(x)$ zu $\mathfrak{I}$ geh"ort.

Es handelt sich um die Frage, was sich "uber die Nenner $a$ dieses maximalen $\mathfrak{M}(\mathfrak{I})$ aussagen l"a\ss t, wenn der urspr"ungliche Integrit"atsbereich $\mathfrak{I}$ ein endlicher war --- d.\,h.\ aus allen ganzzahligen Polynomen endlich vieler Funktionen $f_1(x),\ldots,f_r(x)$ aus $\mathfrak{Z}$ bestand. Dabei sollen f"ur $\mathfrak{Z}$ alle ganzzahligen Polynome oder Potenzreihen in $x_1,\ldots,x_n$ genommen werden, bezw. solche Quotienten, die keine andern als die in den $f(x)$ auftretenden Nenner und deren Potenzen enthalten.

Ich zeige --- unter der im Fall von Potenzreihen oder ihren Quotienten zus"atzlichen Voraussetzung von nur algebraischen Abh"angigkeiten zwischen den $f(x)$ ---, da\ss{} die $a$ sich dann stets so w"ahlen lassen, da\ss{} nur eine endliche Anzahl verschiedener Primzahlen in ihnen aufgeht, diese allerdings im allgemeinen zu jeder Potenz.

Dies letztere ist unmittelbar klar: wenn $a\cdot g(x)$ zu $\mathfrak{I}$ geh"ort, aber $a^{\varrho-1}g(x)^{\varrho}$ f"ur keinen Exponenten $\varrho$, so treten alle Potenzen $a^{\varrho}$ auf (z.\,B. $\mathfrak{I}=[2x]$; $x^{\varrho}=(2x)^{\varrho}/2^{\varrho}$).

Die Frage nach der Beschr"anktheit l"a\ss t sich nur so fassen: L"a\ss t sich ein, im allgemeinen unendliches, Erzeugendensystem von $\mathfrak{M}(\mathfrak{I})$ angeben, derart da\ss{} in den zugeh"origen Nennern die Exponenten der Primzahlen beschr"ankt bleiben? Diese Frage ist --- da es sich um endlich viele Primzahlen handelt --- nach bekannten Schl"ussen\textsuperscript{3)} identisch mit derjenigen nach der Endlichkeit von $\mathfrak{M}(\mathfrak{I})$, also mit dem Hilbert'schen Problem der relativ-ganzen Funktionen in der ganzzahligen Fassung --- eine Frage, die ich mit den angewandten Methoden nicht entscheiden kann\textsuperscript{4)}.

Der Nachweis daf"ur, da\ss{} nur endlich viele verschiedene Primzahlen $p$ in den Nennern $a$ aufgehen, beruht darauf, da\ss{} jedem solchen $p$ mindestens eine Relation modulo $p$ zwischen den $\mathfrak{I}$ erzeugenden $f(x)$ entspricht, eine Relation, die mit ganzzahligen Koeffizienten nicht identisch in $x$ erf"ullt ist. Anders ausgedr"uckt: Bedeutet $\mathfrak{o}$ den ganzzahligen Polynombereich in Unbestimmten $u_1,\ldots,u_r$, so wird $\mathfrak{I}$ verm"oge der Zuordnung $u_i \to f_i(x)$ homomorphes Bild von $\mathfrak{o}$; dieser Homomorphismus wird durch ein Primideal $\mathfrak{a}$ aus $\mathfrak{o}$ vermittelt (d.\,h. $\mathfrak{I}$ wird isomorph dem Restklassenring $\mathfrak{o}/\mathfrak{a}$). Entsprechend wird $\mathfrak{I}_p$ homomorphes Bild von $\mathfrak{o}_p$ --- der Index bedeutet den "Ubergang zu den Restklassen mod.\ $p$, der m"oglich ist mit Ausnahme der endlich vielen in den Nennern von $\mathfrak{I}$ auftretenden Primzahlen. --- Dieser Homomorphismus wird wieder durch ein Primideal $\mathfrak{c}_p$ aus $\mathfrak{o}_p$ vermittelt, das $\mathfrak{a}_p$ umfa\ss t. Dann und nur dann tritt $p$ im Nenner von $\mathfrak{M}(\mathfrak{I})$ auf, wenn $\mathfrak{c}_p$ echtes Oberideal (echter Teiler im Idealsinn) von $\mathfrak{a}_p$ wird. Das kann nur eintreten, wenn sich entweder der Transzendenzgrad von $\mathfrak{I}_p$ gegen"uber $\mathfrak{I}$ erniedrigt oder $\mathfrak{a}_p$ seine Primidealeigenschaft verliert\textsuperscript{5)}. Da\ss{} dies nur f"ur endlich viele $p$ eintritt, folgt daraus, da\ss{} auch modulo $p$ eine algebraische Abh"angigkeit das Verschwinden der Funktionaldeterminante nach sich zieht (\S\,1) und da\ss{} sich weiter $\mathfrak{a}$ als absolutes Primideal ergibt und damit nur modulo endlich vieler $p$ seine Primidealeigenschaft verliert\textsuperscript{6)}.

Es sei noch bemerkt, da\ss{} alle "Uberlegungen mit geringen Modifikationen erhalten bleiben, wenn an Stelle der ganzen rationalen Zahlen ganze Zahlen eines (endlichen) algebraischen Zahlk"orpers und dessen Primideale treten (\S\,4).

% -----------------------------------------------------------
\bigskip
\S~\textbf{1.}
\textbf{Funktionaldeterminanten bei Koeffizientenbereichen der Charakteristik $p$.}

\smallskip

\noindent\textbf{1.}
Es seien $f_1(x_1,\ldots,x_n),\ldots,f_t(x_1,\ldots,x_n)$ rationale Funktionen oder allgemeiner formal definierte Potenzreihen oder Quotienten solcher, mit Koeffizienten aus einem K"orper $P$ der Charakteristik $p$. Die Ableitungen $\partial f_i/\partial x_k$ seien ebenfalls formal definiert, wobei alle "ublichen Rechengesetze erhalten bleiben. Unter $\Omega$ sei ein algebraisch abgeschlossener Erweiterungsk"orper von $P$ verstanden.

Dann gilt ebenso wie bei Charakteristik Null:

\noindent\textbf{Satz.}
\textit{Wenn zwischen den $f(x)$ eine algebraische Abh"angigkeit mit Koeffizienten aus $\Omega$ besteht, so wird der Rang der Funktionalmatrix $(\partial f_i/\partial x_k)$ ($i=1,\ldots,t$; $k=1,\ldots,n$) kleiner als $t$}\textsuperscript{7)}.

\noindent\textit{Beweis.}
Im Polynombereich $\Omega[u_1,\ldots,u_t]$ bedeute $\mathfrak{m}$ das Primideal aller derjenigen Polynome $G(u)$, f"ur die $G(f) = 0$ identisch in $x$. Nach Voraussetzung ist $\mathfrak{m}$ vom Nullideal verschieden; sei $G(u) \neq 0$ (und folglich $\neq {const.}$) ein Polynom niedrigsten Grades aus $\mathfrak{m}$, woraus insbesondere folgt, da\ss{} $G(u)$ unzerlegbar ist. Dann ergibt sich wie "ublich:
\[
\begin{aligned}
\left(\frac{\partial G}{\partial u_1}\right)_{f}\frac{\partial f_1}{\partial x_1}
 +\cdots+
\left(\frac{\partial G}{\partial u_t}\right)_{f}\frac{\partial f_t}{\partial x_1}&=0,\\[-1mm]
&\vdots\\[-1mm]
\left(\frac{\partial G}{\partial u_1}\right)_{f}\frac{\partial f_1}{\partial x_n}
 +\cdots+
\left(\frac{\partial G}{\partial u_t}\right)_{f}\frac{\partial f_t}{\partial x_n}&=0,
\end{aligned}
\]
woraus folgt, da\ss{} entweder der Rang der Funktionalmatrix kleiner als $t$ ist oder aber alle $\left(\partial G/\partial u_i\right)_{f}$ verschwinden. Da aber $G(u)$ als ein Polynom niedrigsten Grades vorausgesetzt war, folgt daraus das Verschwinden aller $\partial G/\partial u_i$.

Somit kommt: $G(u)=H(u_1^p,\ldots,u_t^p)$; und da der Koeffizientenbereich $\Omega$ als algebraisch abgeschlossener, also vollkommener K"orper vorausgesetzt war, weiter: $G(u)=H(u_i^p)=(T(u))^p$; entgegen der Wahl von $G(u)$. --- Der Beweis gilt nat"urlich auch f"ur Charakteristik Null, wobei der letzte Schluss fortf"allt.

\medskip\noindent\textbf{2.} \textbf{Folgerung.}

Es seien $F_1(x_1,\ldots,x_n),\ldots,F_t(x_1,\ldots,x_n)$ ganzzahlige Funktionen (Polynome, Potenzreihen oder Quotienten solcher) mit einer Funktionalmatrix genau vom Rang $t$. Die Primzahl $p$ sei so gew"ahlt, da\ss{} sie nicht in den Nennern der $F(x)$ und nicht in allen $t$-reihigen Determinanten der Funktionalmatrix aufgeht.

Dann bleiben $F_1(x),\ldots,F_t(x)$ auch modulo $p$ algebraisch unabh"angig.

\noindent\textit{Beweis.}
Es bedeute $P$ den Restklassenk"orper nach $p$ und $f_1(x),\ldots,f_t(x)$ die Funktionen aus $P$, die aus $F_1(x),\ldots,F_t(x)$ entstehen, wenn die ganzzahligen Koeffizienten durch ihre Restklassen nach $p$ ersetzt werden. Die $f_i(x)$ sind definiert, da $p$ nicht im Nenner der $F_i(x)$ aufging; ihre Funktionalmatrix $(\partial f_i/\partial x_k)$ entsteht aus $(\partial F_i/\partial x_k)$ ebenfalls durch Ersetzen der ganzzahligen Koeffizienten durch ihre Restklassen nach $p$ --- in $(\partial F_i/\partial x_k)$ treten keine andern Nenner als in den $F(x)$ auf ---. Wegen der Wahl von $p$ beh"alt $(\partial f_i/\partial x_k)$ den Rang $t$; die $f(x)$ sind algebraisch unabh"angig in $\Omega$ und umsomehr in $P$.

% -----------------------------------------------------------
\bigskip
\S~\textbf{2.}
\textbf{Formulierung des Satzes.}

\smallskip

\noindent\textbf{1.}
Wie in der Einleitung bezeichne $\mathfrak{I}$ einen Integrit"atsbereich (Ring ohne Nullteiler) aus ganzzahligen Funktionen in $x_1,\ldots,x_n$ --- Polynomen oder Potenzreihen mit ganzen rationalen Koeffizienten oder Quotienten solcher ---. Im Fall der Potenzreihen kann es sich um formal definierte handeln oder um solche, die in einem gewissen Bereich konvergieren; benutzt wird nur, da\ss{} die auftretenden Funktionen einen Integrit"atsbereich bilden.

Unter $\mathfrak{Q}$ sei der Quotientenk"orper von $\mathfrak{I}$ verstanden; $\mathfrak{Z}$ bedeute einen Integrit"atsbereich zwischen $\mathfrak{I}$ und $\mathfrak{Q}$. Wie in der Einleitung hei\ss e $\mathfrak{M}$ \textit{maximal} in $\mathfrak{Z}$, wenn aus der Zugeh"origkeit von $a \cdot g(x)$ --- $a$ eine ganze Zahl $\neq 0$, $g(x)$ in $\mathfrak{Z}$ --- stets diejenige von $g(x)$ folgt. Zu $\mathfrak{I}$ existiert eindeutig der durch $\mathfrak{I}$ in $\mathfrak{Z}$ erzeugte Maximalbereich $\mathfrak{M}(\mathfrak{I})$ als Durchschnitt aller $\mathfrak{I}$ umfassenden Maximalbereiche in $\mathfrak{Z}$ --- denn $\mathfrak{Z}$ selbst ist ein solcher, und mit beliebig vielen ist auch der Durchschnitt maximal und $\mathfrak{I}$ umfassend. $\mathfrak{M}(\mathfrak{I})$ besteht aus allen Funktionen $g(x)$ aus $\mathfrak{Z}$, f"ur die es mindestens ein $a \neq 0$ gibt, so da\ss{} $a \cdot g(x)$ in $\mathfrak{I}$ liegt.

Denn diese $g(x)$ bilden einen $\mathfrak{I}$ umfassenden Maximalbereich $\mathfrak{N}$ in $\mathfrak{Z}$, da f"ur $b \cdot h(x)$ in $\mathfrak{I}$ mit $b \neq 0$ und $h(x)$ in $\mathfrak{Z}$ folgt, da\ss{} $ab \cdot h(x)$ in $\mathfrak{I}$ liegt und $ab \neq 0$ ist, also $h(x)$ in $\mathfrak{N}$ liegt. Jeder $\mathfrak{I}$ umfassende Maximalbereich $M$ in $\mathfrak{Z}$ aber umfa\ss t $\mathfrak{N}$; denn liegt $a \cdot g(x)$ in $\mathfrak{I}$, so liegt $a \cdot g(x)$ in $\mathfrak{M}$, da $\mathfrak{I}$ in $\mathfrak{M}$ liegt; also $g(x)$ in $\mathfrak{M}$.

\medskip\noindent\textbf{2.}
Sei jetzt $\mathfrak{I}$ als endlicher Integrit"atsbereich mit Einheitselement vorausgesetzt, mit $f_1(x),\ldots,f_r(x)$ als Integrit"atsbasis; d.\,h.\ $\mathfrak{I}$ bestehe aus allen ganzzahligen Polynomen der $f_i(x)$. Weiter bestehe $\mathfrak{Z}$ aus allen ganzzahligen Polynomen oder Potenzreihen aus $x$, bezw.\ aus solchen Quotienten, die keine andern als die in den endlich vielen $f_i(x)$ auftretenden Nenner enthalten; diese nat"urlich zu jeder Potenz.

Ferner sei angenommen, da\ss{} in mindestens einem $f_i(x)$ --- im allgemeinen in allen --- die $x$ wirklich auftreten, da sonst einfach $\mathfrak{I}$ und $\mathfrak{Z}$ aus allen ganzzahligen Polynomen einer gebrochenen Zahl $1/c$ bestehen und nichts zu beweisen ist.

Handelt es sich um Polynome oder rationale Funktionen, so wird also der Quotientenk"orper $\mathfrak{Q}$ vom Transzendenzgrad $t$ mit $1 \leq t \leq r$ in bezug auf den K"orper $K$ der rationalen Zahlen; und wenn etwa $f_1(x),\ldots,f_t(x)$ ein irreduzibles System (Anmerkung\textsuperscript{5)}) bilden, so wird $\mathfrak{Q}$ endliche algebraische Erweiterung der rein transzendenten $K(f_1(x),\ldots,f_t(x))$. Zugleich hat die Funktionalmatrix von $f_1(x),\ldots,f_r(x)$ und ebenso die von $f_1(x),\ldots,f_t(x)$ genau den Rang $t$\textsuperscript{8)}.

Damit auch im Fall von Potenzreihen oder ihren Quotienten diese Struktur von $\mathfrak{Q}$ erhalten bleibt, ist eine Zusatzvoraussetzung zu machen: Ist $t$ der Rang der Funktionalmatrix der $f_i(x)$ und etwa zugleich der Rang der Funktionalmatrix von $f_1(x),\ldots,f_t(x)$ --- sodass diese nach \S\,1, 1.\ ein irreduzibles System bilden ---, so soll wieder $\mathfrak{Q}$ endliche algebraische Erweiterung der rein transzendenten $K(f_1(x),\ldots,f_t(x))$ sein; d.\,h.\ $f_{t+1}(x),\ldots,f_r(x)$ sollen algebraisch in bezug auf $K(f_1(x),\ldots,f_t(x))$ sein\textsuperscript{9)}. Dabei wird wieder $t \geq 1$, da --- Charakteristik Null! --- nicht alle $\partial f_i/\partial x_k$ verschwinden.

Zugleich folgt die Bedingung f"ur jedes System von $t$ Funktionen der Integrit"atsbasis, deren Funktionalmatrix genau vom Rang $t$ ist; denn diese bilden ein irreduzibles System, von dem also die "ubrigen algebraisch abh"angen.

